\section{基于A-D3QN的混合路径规划算法}
\label{sec:ch6_ad3qn}

为兼顾全局最优性与局部实时性，本文采用 A* 与 D3QN 协同的混合规划策略：A* 提供全局可行骨架，D3QN 在局部动态不确定区域内执行细粒度动作决策。
该策略在系统层体现为“全局先验生成-局部策略执行-反馈校正”的协同闭环。

\begin{figure}[H]
	\centering
	\resizebox{0.86\textwidth}{!}{%
	\begin{tikzpicture}[font=\small, >=Stealth]
		\node[draw, rounded corners, align=center, minimum width=3.2cm, minimum height=1.0cm, fill=gray!10] (in) at (0,0) {环境状态输入\\$s_t$};

		\node[draw, rounded corners, align=center, minimum width=3.5cm, minimum height=1.0cm, fill=blue!10] (a1) at (-4.4,-1.8) {全局地图建模};
		\node[draw, rounded corners, align=center, minimum width=3.5cm, minimum height=1.0cm, fill=blue!10] (a2) at (-4.4,-3.4) {改进 A* 搜索};
		\node[draw, rounded corners, align=center, minimum width=3.5cm, minimum height=1.0cm, fill=blue!10] (a3) at (-4.4,-5.0) {路径平滑与约束校核};

		\node[draw, rounded corners, align=center, minimum width=3.8cm, minimum height=1.0cm, fill=green!10] (d1) at (4.4,-1.8) {局部观测编码};
		\node[draw, rounded corners, align=center, minimum width=3.8cm, minimum height=1.0cm, fill=green!10] (d2) at (4.4,-3.4) {Dueling Double Q 网络};
		\node[draw, rounded corners, align=center, minimum width=3.8cm, minimum height=1.0cm, fill=green!10] (d3) at (4.4,-5.0) {动作价值评估};

		\node[draw, diamond, aspect=2.2, align=center, inner sep=1.2pt, fill=orange!15] (sw) at (0,-3.4) {切换判据\\$\mathcal{S}_t$};
		\node[draw, rounded corners, align=center, minimum width=4.4cm, minimum height=1.0cm, fill=purple!10] (out) at (0,-6.8) {动作输出与执行反馈\\$a_t,\ e_t$};

		\draw[->, thick] (in) -- (a1);
		\draw[->, thick] (in) -- (d1);
		\draw[->, thick] (a1) -- (a2);
		\draw[->, thick] (a2) -- (a3);
		\draw[->, thick] (d1) -- (d2);
		\draw[->, thick] (d2) -- (d3);
		\draw[->, thick] (a2.east) -- (sw.west);
		\draw[->, thick] (d2.west) -- (sw.east);
		\draw[->, thick] (a3.east) .. controls (-2.4,-5.0) and (-1.6,-6.4) .. (out.west);
		\draw[->, thick] (d3.west) .. controls (2.4,-5.0) and (1.6,-6.4) .. (out.east);
		\draw[->, thick] (out.north) .. controls (0.0,-5.8) and (0.0,-4.5) .. node[right]{在线校正} (sw.south);
	\end{tikzpicture}%
	}
	\caption{A*-D3QN-Opt 混合路径规划算法架构}
	\label{fig:ch6_ad3qn_arch}
\end{figure}

图\ref{fig:ch6_ad3qn_arch}展示了全局搜索、局部强化学习与反馈校正三部分的协同关系。
其中 A* 负责生成可达性先验，D3QN 负责局部动作优化，反馈模块依据执行偏差触发重规划或参数更新，
从而兼顾路径全局最优趋势与局部实时响应能力。

\subsection{传统A*算法改进}
\label{subsec:ch6_astar_improve}

项目 A* 前端采用 8 邻域搜索，并在代价函数中引入能耗项：
\begin{equation}
g(n)=\sum_{k=1}^{n}\left(c_{dist,k}+0.1\,c_{eng,k}\right),\qquad
f(n)=g(n)+h(n),
\label{eq:ch6_astar_cost}
\end{equation}
其中 $h(n)$ 为欧氏距离启发项。
当起点或终点落在障碍栅格时，算法通过最近可行栅格搜索进行鲁棒修正，提升任务可执行性。

为抑制网格路径折线震荡，采用梯度下降平滑：
\begin{equation}
\mathbf{p}_{i}^{t+1}=\mathbf{p}_{i}^{t}
+w_d\left(\mathbf{p}_{i}^{0}-\mathbf{p}_{i}^{t}\right)
+w_s\left(\mathbf{p}_{i-1}^{t}+\mathbf{p}_{i+1}^{t}-2\mathbf{p}_{i}^{t}\right),
\label{eq:ch6_path_smooth}
\end{equation}
其中 $w_d=0.5$、$w_s=0.1$。

在端到端实验中，平滑前后对比如下：路径长度由 $1289.19\,\mathrm{m}$ 降至 $1280.99\,\mathrm{m}$，累计转向量由 $79.33$ 降至 $50.96$，说明改进 A* 在几乎不增加计算代价的前提下显著提升可执行性。

\subsection{D3QN深度强化学习算法}
\label{subsec:ch6_d3qn}

D3QN 网络采用“卷积编码 + Dueling 结构 + Double DQN 目标”。
输入为深度帧堆叠 $\mathbf{S}\in \mathbb{R}^{4\times80\times64}$ 与 6 维附加特征（目标方向/距离、A*方向/距离、障碍方向/距离）。

Q 值分解为
\begin{equation}
Q(s,a;\theta)=V(s;\theta_v)+\left(A(s,a;\theta_a)-\frac{1}{|\mathcal{A}|}\sum_{a'}A(s,a';\theta_a)\right).
\label{eq:ch6_dueling_q}
\end{equation}

Double DQN 目标值写为
\begin{equation}
y_t=r_t+\gamma(1-d_t)\,Q_{\theta^-}\!\left(s_{t+1},
\arg\max_{a}Q_{\theta}(s_{t+1},a)\right),
\label{eq:ch6_ddqn_target}
\end{equation}
损失函数为
\begin{equation}
\mathcal{L}=\mathbb{E}\left[w_i\left(Q_{\theta}(s_t,a_t)-y_t\right)^2\right],
\label{eq:ch6_d3qn_loss}
\end{equation}
其中 $w_i$ 为优先经验回放权重。

工程参数与实现细节如下：
\begin{myitemize}
	\item 网络参数量：$1{,}146{,}762$；
	\item 学习率：$1\times10^{-4}$，折扣因子 $\gamma=0.99$；
	\item 回放池容量：$100{,}000$，批量大小：$32$；
	\item 探索率：$\epsilon:1.0\rightarrow0.01$（指数衰减）；
	\item 目标网络软更新：$\tau=0.001$（每4步更新）。
\end{myitemize}

在 CPU 条件下，D3QN 单步前向推理均值为 $1.21\,\mathrm{ms}$（$P95=1.96\,\mathrm{ms}$），估算推理吞吐约 $828.78\,\mathrm{FPS}$，满足局部决策实时性要求。

训练奖励采用多项构成：
\begin{equation}
r_t=r_{base}+5\Delta d_t-r_{step}+r_{prox}+r_{goal},
\label{eq:ch6_reward}
\end{equation}
其中 $\Delta d_t$ 为目标距离改变量。
权重敏感性分析表明，动力学惩罚权重最优点为 $w_{dyn}=0.45$。

\subsection{混合算法融合策略}
\label{subsec:ch6_fusion}

融合策略由分层智能体实现：先计算全局路径，再根据当前位置选取前视局部目标，并将 A* 先验作为附加特征输入 D3QN 网络。
关键切换参数为：前视点数 $\mathrm{look\_ahead}=3$、路径偏离重规划阈值 $2.0\,\mathrm{m}$、到点阈值 $1.0\,\mathrm{m}$。

混合决策可写为
\begin{equation}
\mathbf{u}_k=
\begin{cases}
\pi_{D3QN}\!\left(\mathbf{S}_k,\mathbf{g}^{A*}_k\right), & \delta_k\le \tau_{replan},\\
\pi_{A*}\!\left(\mathcal{M}_{DT},\mathbf{x}_k,\mathbf{x}_g\right), & \delta_k> \tau_{replan},
\end{cases}
\label{eq:ch6_switch}
\end{equation}
其中 $\delta_k$ 为当前位置到全局路径最近点距离。

\begin{algorithm}[H]
\caption{A-D3QN 混合路径规划}
\label{alg:ch6_ad3qn}
\begin{algorithmic}[1]
\REQUIRE 数字孪生地图 $\mathcal{M}_{DT}$，起点 $\mathbf{x}_s$，目标 $\mathbf{x}_g$
\ENSURE 控制序列 $\{\mathbf{u}_k\}$
\STATE 使用 A* 生成全局路径 $\mathcal{P}_g$
\STATE 对 $\mathcal{P}_g$ 执行平滑，初始化前视索引
\FOR{$k=0,1,\dots$}
\STATE 计算偏离距离 $\delta_k$ 与局部目标点
\IF{$\delta_k>\tau_{replan}$}
\STATE 触发 A* 重规划并更新 $\mathcal{P}_g$
\ENDIF
\STATE 构造局部状态 $\mathbf{S}_k$ 与 A* 附加特征 $\mathbf{g}^{A*}_k$
\STATE 由 D3QN 输出动作 $\mathbf{u}_k$
\STATE 执行动力学预测并回传滑移/沉陷残差
\IF{到达终点}
\STATE \textbf{break}
\ENDIF
\ENDFOR
\end{algorithmic}
\end{algorithm}

该融合机制使系统同时具备全局可达性保障与局部动态响应能力，为第\ref{sec:ch6_validation}节实验验证提供算法基础。
